\section{Conclusion and Future Work}

We introduced and evaluated the Evolutionary Neuroplastic Agent (ENA), an algorithmic framework fundamentally designed to combat catastrophic interference over shifting environments by maintaining a highly responsive, continuously updated archive of sub-specialists governed competitively through heuristic \textit{Trust} and \textit{Plasticity} limits. Extensive evaluation across changing physical gravities established that deploying Fuzzy-logic metric derivations substantially minimized the observed Stability Gap during domain adjustments (down to $-2.13$) when juxtaposed against deterministic quadratic ratio limits. The mechanism allows an algorithmic bridging between pure Quality-Diversity maps and direct continuous reinforcement agents.

Crucially, experimentation unveiled the direct limitations dictated by initialization variance on populations. Constraining learning vectors to limited search spaces can produce unpredictable isolated performance degradation (represented by instances diving toward $0.80\%$ capability), heavily supporting the hypothesis that robust adaptation demands substantial initialized genetic diversity.

Future work will expand the analysis beyond the preliminary 1D kinematic CartPole benchmark by scaling generation counts and population pools sequentially into deeper continuous domains. By mathematically eroding the initialization bottleneck, we plan to directly benchmark ENA's overarching rapid-learning mechanisms continuously against established baseline DRL frameworks (such as PPO or Deep Q-Networks) to validate absolute operational supremacy.
